% =========================================================================
% $I$ Introduction (lane C).  Written 2026-09-30 for the merged two-lane note.
% =========================================================================
\section{Introduction}\label{sec:intro}

\emph{The object.} The reference paper \cite{refpaper} studies a remote loop in
which the sensor is \emph{virtual}: it is described by an information matrix $S$
injected per step, and the designer is restricted to what the task actually
reads, i.e.\ $S$ must lie in the cone
$\mathcal K_F=\{S\succeq0:\operatorname{range}(S)\subseteq
\operatorname{range}(F^{\top})\}$. The paper plots rate--cost frontiers for a
published four-state example, parameterises the cone (its Corollary~1), and its
Theorem~3 closes the loop: at \emph{every} cost level $D$ the restricted design
attains the unrestricted optimum, so the task restriction costs nothing.

\emph{What we find.} We reproduce the published frontier over the whole cone ---
including the rank-deficient boundary where channels are switched off --- and
then report what the published window can and cannot show.
\begin{enumerate}\itemsep2pt
\item The frontier is \emph{not} reproduced by any isotropic-noise family: the
      best member of each family lies $1.11$--$16.4$ cost units above it, while
      optima over the full cone agree with it to within $0.07$ units, inside the
      digitisation band (\S\ref{sec:repro}).
\item The admissibility wall is a \emph{horizontal asymptote} of each frontier,
      not a point on it: the cost approaches $D_{\min}(V)$ from above and never
      reaches it, so the published figure cannot display the wall it names. The
      wall is a functional of the subspace alone,
      $D_{\min}(V)=j_c+\Phi_0(V)$, and the interval $(j_c,D_{\min}(V)]$ is a
      \emph{complete infeasible region}: inside it no admissible design meets the
      task (\S\ref{sec:infeasible}).
\item Consequently the phrase ``for every cost level $D$'' in Theorem~3 needs the
      hypothesis its own proof already uses, that the stage cost sees the state
      only through the task variable, $\operatorname{supp}(Q)\subseteq
      \operatorname{range}(F^{\top})$. Without it the theorem fails not at a point
      but on an interval of length exactly $\Phi_0(V)$; the repaired statement
      (Theorem~3$'$) is free and gives both $\Delta\equiv0$ and $\Phi_0(V)=0$
      (\S\ref{sec:infeasible}).
\item The wall does not order the frontiers: among $120$ random rank-one designs,
      $111$ pairs invert the wall ordering at a fixed rate, so the wall is not a
      design criterion (\S\ref{sec:nofactor}).
\item On the constructive side, a frozen prior reduces the trade-off to one convex
      program on the information form, whose solution is reverse water-filling on
      the spectrum of a task-weighted Gramian $\Gamma$. This supplies a
      closed-form value at the \emph{wall prior}, not the previously
      unspecified constant of the wall-approach law --- the frontier-wide
      form that would have delivered that constant needs a monotonicity that
      is \emph{false} (a closed-form rank-one counterexample and $115$ reachable
      in-cone counterexamples), and we withdraw it --- together with a closed-form
      sensor that
      matches the best previously known designs at ten of eleven anchors to within
      $1.4\%$ and reproduces the rank-deficiency of the numerically optimal designs
      through channel switch-on rates (\S\ref{sec:waterfill}).
\item On the floor side, choosing the sensing subspace itself is a
      lower-dimensional problem with a usable rule and a certification status we
      state precisely: the rule attains
      $\Phi_0=12.222448/0.791828/0.190335$ at rank $1/2/3$ against best known
      $10.169224/0.735143/0.158701$, i.e.\ $1.20/1.08/1.20\times$ the best known
      floor (\S\ref{sec:floor}).
\item On the rate side we characterise the restricted sensor beyond the published
      window. With the two excesses $x:=D-D_{\min}(V)>0$ (cost axis) and
      $\Delta I:=I-R_{\exp}$ (rate axis) --- different origins, never mixed --- the
      ladder law
      $\Delta I=\frac{\mathrm{rank}C}{2}\log_2\frac1x+b_{\rm pred}+O(x)$ has an
      exactly computable intercept and asymptotic slope
      $\gamma_\infty=2\ln2/r$, which we verify on two independently
      \emph{measurable} legs across five plants and ranks $1$--$4$: the rate leg
      $\sigma:=d\Delta I/d\log_2s\to r/2$ and the cost leg
      $\alpha:=-d\ln x/d\ln s\to1$ (\S\ref{sec:rate}). The slope itself is
      reverse-water-filling arithmetic and is \emph{not} claimed as a discovery;
      what we claim is its domain, the finite-rate inequality
      $\gamma_r/\gamma_1>1/r$, the intercept, and the rank structure.
      The same section records where the water-filling picture fails: the optimal
      support is \emph{not} diagonal in the task-weight basis, and forcing it
      there costs $5.0\times10^{-2}$/$9.5\times10^{-2}$ bit at rank one.
\item Finally we reduce the one question the published line leaves open ---
      whether admissibility can ever cost rate --- to a statement about the
      \emph{uniqueness of the optimal support}, which random probes of the
      optimal face cannot decide by construction, and we report its direct
      numerical decision with a pre-registered criterion
      (\S\ref{sec:rate}, Appendix~\ref{app:repro}).
\end{enumerate}

\emph{Positioning.} This is a comment-and-extension note, not a competing design.
We keep the reference paper's model, symbols and cone; we reproduce its
Figure~1 over a larger admissible set than it displays; we correct one hypothesis
in its Theorem~3; and we extend the picture in two directions it does not take
(the subspace-selection floor side, and the ladder law with its rank structure on
the rate side). Where its own moving parts are ours we say so: its Step~3
realizability test is the same equality-type reachability condition we use, and
its determinant monotonicity is a statement about one PSD variable under a Loewner
bound, not the $\det\Gamma$ monotonicity: that one is \emph{false} (closed-form
rank-one counterexample; $115$ reachable in-cone counterexamples), which is why
the frontier-wide bound resting on it is withdrawn in \S\ref{sec:waterfill}
(\S\ref{sec:related}). Relative to the closest prior on restricting sensing
\cite{cuvelier2021side} we add the cone language and the \emph{length} of the
infeasible interval, not the act of restricting sensing; relative to the
nonanticipative rate-distortion line, we price a chosen sensing subspace against a
control cost, and we claim no converse of our own.

\emph{Honesty.} Everything quantitative is for the published four-state example.
All rate-side numbers are constructive upper bounds obtained by local descent,
with a lower-bound certificate only where explicitly said to be proved. The
floor-side rule is ``best known, certified global on the floor side''. Every
statement we had to withdraw, and every limitation we know of, is collected in
\S\ref{sec:qual}; the claim-to-script-to-log map is Appendix~\ref{app:repro}.
